%File: PRISM-Main-2026.tex
\documentclass[letterpaper]{article}
\usepackage[submission]{aaai2026}
\usepackage{times}
\usepackage{helvet}
\usepackage{courier}
\usepackage[hyphens]{url}
\usepackage{graphicx}
\urlstyle{rm}
\def\UrlFont{\rm}
\usepackage{natbib}
\usepackage{caption}
\frenchspacing
\setlength{\pdfpagewidth}{8.5in}
\setlength{\pdfpageheight}{11in}

\usepackage{algorithm}
\usepackage{algorithmic}

\usepackage{newfloat}
\usepackage{listings}
\DeclareCaptionStyle{ruled}{labelfont=normalfont,labelsep=colon,strut=off}
\lstset{%
    basicstyle={\footnotesize\ttfamily},
    numbers=left,numberstyle=\footnotesize,xleftmargin=2em,
    aboveskip=0pt,belowskip=0pt,
    showstringspaces=false,tabsize=2,breaklines=true}
\floatstyle{ruled}
\newfloat{listing}{tb}{lst}{}
\floatname{listing}{Listing}

\pdfinfo{
/TemplateVersion (2026.1)
}

\setcounter{secnumdepth}{0}

% ============================================================================
% TITLE AND AUTHORS
% ============================================================================

\title{PRISM: Paradigm-Guided Reasoning with Iterative Solver Memory for Constraint Satisfaction}

\author{
Anonymous AAAI Submission
}

% ============================================================================
% ABSTRACT
% ============================================================================

\begin{document}

\maketitle

\begin{abstract}
Constraint Satisfaction Problems (CSPs) are fundamental to numerous real-world applications spanning task planning, resource allocation, and scheduling optimization. While Large Language Models (LLMs) have demonstrated remarkable capabilities in natural language understanding and mathematical reasoning, they face fundamental challenges in structured constraint reasoning. Recent systematic evaluations reveal a "curse of complexity" phenomenon: as puzzle scale increases, all tested models experience sharp accuracy degradation, approaching random performance when the Z3 solver encounters more than 20 conflicts. This work proposes \textbf{PRISM} (Paradigm-guided Reasoning with Iterative Solver Memory), a novel neuro-symbolic framework that addresses two critical limitations of existing approaches: (1) the inefficiency of repair mechanisms that lack semantic feedback, and (2) the absence of experience accumulation across problem instances.

PRISM employs a two-stage architecture. In the \textit{offline phase}, the system distills structured solving paradigms from successful trajectories via Key Decision Point (KDP) identification, cross-trajectory clustering, and Z3-based triple verification (soundness, effect, trigger precision). In the \textit{online phase}, retrieved paradigms guide LLM inference while a typed repair trajectory memory employs Jaccard-based stagnation detection and cosine-based loop detection to trigger a four-level strategy escalation protocol. Crucially, paradigms encode formal Z3 pre/post-conditions enabling mechanical verification—a fundamental distinction from prior memory-based approaches relying on natural language heuristics.

Evaluation on the ZebraLogic benchmark across six complexity scales (3×5 to 6×6) demonstrates significant improvements over seven baselines, with the largest gains on the most difficult 6×6 puzzles, validating our core thesis that experience accumulation becomes increasingly valuable for harder problems. The approach achieves a 40:1 compression ratio (1500+ training trajectories to ~35 verified paradigms) and reduces average repair rounds by 40-50\% compared to neural-symbolic baselines.
\end{abstract}

% ============================================================================
% INTRODUCTION
% ============================================================================

\section{Introduction}

Constraint Satisfaction Problems (CSPs) form the backbone of numerous real-world applications including task planning, resource allocation, and combinatorial optimization. Formally, a CSP is defined as a tuple $\mathcal{P} = (\mathbf{X}, \mathbf{D}, \mathbf{C})$ where $\mathbf{X}$ is a set of variables, $\mathbf{D}$ specifies finite domains for each variable, and $\mathcal{C}$ is a set of constraints. A solution is a complete assignment of variables that simultaneously satisfies all constraints.

Recent advances in Large Language Models have demonstrated impressive capabilities across diverse reasoning tasks. However, LLMs struggle with structured constraint reasoning—a phenomenon systematically documented in recent benchmarks. The ZebraLogic benchmark (Lin et al., 2025) reveals a critical limitation: as puzzle complexity scales from 3×5 to 6×6 logical grids, all evaluated models exhibit sharp accuracy degradation. When the Z3 constraint solver encounters more than 20 conflicts, model performance approaches random baselines, and computational scaling techniques (Best-of-N sampling, self-verification) yield diminishing returns.

This performance cliff is not merely a scale-dependent phenomenon but stems from a fundamental architectural limitation: LLMs lack systematic strategies for navigating complex constraint spaces. Existing neural-symbolic approaches (Pan et al., 2023; Olausson et al., 2023) adopt a three-stage translate-execute-repair pipeline where LLMs convert natural language constraints to formal representations, solvers perform deduction, and LLMs attempt repairs upon failure. This paradigm has achieved improvements over pure language model reasoning but exhibits two critical deficiencies:

\textbf{(1) Repair Inefficiency:} When solvers return UNSAT, repair mechanisms receive only coarse feedback (e.g., raw error strings or generic unsatisfiability statements) lacking semantic information about constraint conflicts. LLMs thus iteratively generate near-identical erroneous repairs—a phenomenon documented as "stagnation" in recent work. Critically, the system cannot distinguish whether contradictions arise from "unsolvable formulations" or "over-constrained translations," nor can it detect repair cycles.

\textbf{(2) Absence of Experience Accumulation:} Each puzzle is solved independently; successful translation patterns, error diagnoses, and repair strategies accumulated on one problem are completely discarded when switching to the next. Identical constraint translation errors recur across different problem instances, necessitating redundant repair iterations.

Agent memory methods (Shinn et al., 2023; Zhao et al., 2024; Ouyang et al., 2025) explore cross-task experience accumulation through trajectory comparison and heuristic extraction. However, existing memory frameworks face inherent limitations when applied to CSP reasoning: (a) memory units are unverified natural language prose lacking formal correctness guarantees, and (b) task-level memory granularity mismatches the operation-level semantic patterns required for constraint reasoning.

\textbf{Key Insight:} CSP solving exhibits a unique property absent in general agent tasks: every reasoning step is mechanically verifiable via SMT solvers like Z3. This property fundamentally transforms experience accumulation: paradigms can receive automated correctness certification before storage (rather than relying on subjective LLM evaluation), and repair failures generate UNSAT cores—minimal constraint subsets responsible for contradictions—providing precise semantic localization.

We introduce \textbf{PRISM} (Paradigm-guided Reasoning with Iterative Solver Memory), a memory-augmented neuro-symbolic framework designed specifically for CSP solving. PRISM operates via two decoupled phases:

\textbf{Offline Phase (one-time execution):} Given a corpus of training puzzles, the system:
\begin{enumerate}
    \item Collects successful solving trajectories via LLM+Z3 interaction
    \item Identifies Key Decision Points (KDPs)—steps where domain reductions are significant or inference types are non-trivial
    \item Clusters KDPs via cross-trajectory similarity on constraint-type feature vectors
    \item Abstracts reusable solving paradigms via LLM synthesis
    \item Filters candidates through triple Z3 verification (soundness, effect, trigger precision)
    \item Persists verified paradigms in a Paradigm Library
\end{enumerate}

\textbf{Online Phase (per-puzzle execution):} For each target puzzle, PRISM:
\begin{enumerate}
    \item Extracts constraint-type signatures from current solver states
    \item Retrieves relevant paradigms via two-layer filtering (type matching + semantic verification)
    \item Injects paradigms as structured hints to guide LLM inference
    \item Performs real-time Z3 verification of each reasoning step
    \item Upon UNSAT outcomes, extracts UNSAT cores for causal analysis
    \item Records typed repair attempts in a Repair Trajectory Memory
    \item Detects stagnation (Jaccard similarity of UNSAT cores) and loops (cosine similarity of repair actions)
    \item Triggers four-level strategy escalation (L1: switch repair target; L2: switch repair type; L3: revert to recent SAT checkpoint; L4: full retranslation)
    \item Writes back successful repairs to augment the Paradigm Library
\end{enumerate}

\textbf{Main Contributions:}

\begin{enumerate}
    \item \textbf{Formally Verifiable Paradigm Representation}: We introduce the first paradigm formalism for CSP reasoning with mechanical Z3 verification. Paradigms encode trigger conditions, inference operations, and Z3 pre/post-conditions, enabling quality certification impossible in prior natural-language memory approaches.

    \item \textbf{Automated Paradigm Distillation Pipeline}: We design a fully unsupervised offline pipeline that compresses 1500+ training trajectories into ~35 high-confidence paradigms (40:1 ratio) via KDP identification, clustering, LLM abstraction, and formal verification.

    \item \textbf{Typed Repair Trajectory Memory with Adaptive Escalation}: We propose a structured memory system that models constraint repair as learnable trajectories, incorporating Jaccard-based stagnation detection, cosine-based loop detection, and a four-level strategy escalation protocol that mechanically breaks repair stagnation.

    \item \textbf{Paradigm Transferability Evaluation Framework}: We establish "paradigm transfer rate" as a novel evaluation dimension, quantifying cross-scale (L1) and cross-problem (L2) generalization, providing methodological foundations for neuro-symbolic experience accumulation research.
\end{enumerate}

\textbf{Experimental Results:} Evaluation on ZebraLogic (6 complexity scales, 3×5 to 6×6) demonstrates consistent improvements over seven baselines including pure CoT, logic-specific methods (Logic-LM, ExpeL), and neural-symbolic baselines. The largest gains appear on the most difficult 6×6 puzzles, validating our hypothesis that experience accumulation provides exponentially increasing value for harder problems.

The remainder of this paper is organized as follows. Section~2 reviews related work on neural-symbolic reasoning and agent memory. Section~3 presents the PRISM methodology in detail. Section~4 describes experimental setup and baselines. Section~5 reports results and analysis. Section~6 concludes with implications and future directions.

% ============================================================================
% RELATED WORK
% ============================================================================

\section{Related Work}

\subsection{Neural-Symbolic Constraint Reasoning}

Early neural-symbolic approaches (Mao et al., 2019; Yi et al., 2019) combined perception networks with symbolic programs. Recent work focuses on LLM-based constraint reasoning. Pan et al.~(2023) propose Logic-LM, combining language models with symbolic executors for mathematical reasoning. Their translate-execute-repair pipeline forms the foundation for subsequent neuro-symbolic CSP work. Olausson et al.~(2023) employ Z3 for explicit constraint verification, while Ye et al.~(2023) integrate formal methods into language model reasoning pipelines.

However, existing approaches exhibit two critical limitations: (1) repair mechanisms receive only coarse error signals, leading to stagnation; (2) each problem is solved independently without cross-problem experience transfer. PRISM addresses both through formally-grounded paradigm libraries and typed repair memory.

\subsection{Agent Memory and Experience Accumulation}

Memory-augmented agents (Shinn et al., 2023; Zhao et al., 2024; Ouyang et al., 2025) accumulate experiences across tasks through trajectory comparison and heuristic extraction. These methods demonstrate effective cross-task transfer in web navigation, embodied tasks, and software engineering.

Critically, existing memory frameworks have two limitations when applied to CSP reasoning: (a) memory units are unverified natural-language heuristics, and (b) task-level granularity mismatches the operation-level semantic patterns required. PRISM overcomes these through Z3-verified paradigms and operation-level repair records.

\subsection{Formal Methods and Constraint Solving}

Classical CSP solvers (Dechter, 2003; Russell \& Norvig, 2021) employ systematic search and constraint propagation. SMT solvers like Z3 (de Moura \& Bjørner, 2008) provide mechanical verification of logical formulas. Recent work increasingly leverages formal methods: Talmor et al.~(2019) use automated theorem provers for reasoning, while Sinha et al.~(2021) combine neural networks with formal verification.

PRISM's key innovation is leveraging Z3 as a \textit{quality oracle} for paradigm verification, not merely as a solver. This enables mechanical certification of experience without human annotation.

% ============================================================================
% METHODOLOGY
% ============================================================================

\section{Methodology}

\subsection{Problem Definition and Notation}

A Constraint Satisfaction Problem is formally defined as $\mathcal{P} = (\mathbf{X}, \mathbf{D}, \mathbf{C})$ where:
\begin{itemize}
    \item $\mathbf{X} = \{x_1, x_2, \ldots, x_n\}$ is a set of variables
    \item $\mathbf{D} = \{D_1, D_2, \ldots, D_n\}$ specifies finite domains ($x_i \in D_i$)
    \item $\mathcal{C} = \{c_1, c_2, \ldots, c_m\}$ is a set of constraints
\end{itemize}

A solution is a complete assignment $\theta: \mathbf{X} \to \bigcup_i D_i$ such that $\theta(x_i) \in D_i$ and $\forall c_j \in \mathcal{C}: c_j(\theta)$ holds.

\textbf{Solving Trajectory:} For puzzle instance $\mathcal{P}$, a solving trajectory is a temporal sequence of reasoning states:
$$\mathcal{T} = \langle (S_0, a_1, S_1), (S_1, a_2, S_2), \ldots, (S_{T-1}, a_T, S_T) \rangle$$

where $S_t = (C_t, \delta_t)$ is the solver state at step $t$, with $C_t$ being formalized constraints and $\delta_t$ being domain assignments. Action $a_t$ is LLM-generated inference (natural language statement + Z3 formalization). Trajectory $\mathcal{T}$ is \textit{successful} if $S_T$ satisfies all constraints with Z3 returning SAT and assignment matching ground truth.

\textbf{Solving Paradigm:} We define a solving paradigm as a five-tuple:
$$P = (\mathcal{Q}, \mathcal{O}, \phi_{\text{pre}}, \phi_{\text{post}}, \mathcal{S})$$

where:
\begin{itemize}
    \item $\mathcal{Q}$ specifies trigger conditions (constraint types and solver state features)
    \item $\mathcal{O}$ is the inference operation (parameterized template, e.g., $\texttt{pos}(X) = k + d$)
    \item $\phi_{\text{pre}}$ is the Z3-formalized precondition assertion
    \item $\phi_{\text{post}}$ is the Z3-formalized postcondition assertion
    \item $\mathcal{S}$ documents applicable scope (puzzle types and scales)
\end{itemize}

Critically, $\phi_{\text{pre}}$ and $\phi_{\text{post}}$ have formal semantics and are mechanically verifiable by Z3—a fundamental distinction from prior natural-language memory work.

\subsection{PRISM Framework Overview}

PRISM operates via two decoupled stages: offline distillation and online inference.

\textbf{Offline Phase:} Given training puzzles $\{\mathcal{P}_1, \ldots, \mathcal{P}_N\}$:
\begin{enumerate}
    \item Collect successful solving trajectories via LLM+Z3
    \item Identify Key Decision Points (domain reductions ≥2 or non-trivial step types)
    \item Cluster KDPs via complete-linkage agglomerative clustering on constraint-type features
    \item Abstract paradigms via LLM synthesis from cluster representatives
    \item Filter candidates through triple Z3 verification (soundness ≥0.90, effect ≥0.80, precision ≥0.80)
    \item Persist verified paradigms in Paradigm Library $\mathcal{L}$
\end{enumerate}

\textbf{Online Phase:} For each puzzle $\mathcal{P}^*$:
\begin{enumerate}
    \item Extract constraint-type signature $\sigma(S_t) = \text{sorted}(\{type(c) \mid c \in C_t\})$
    \item Retrieve paradigms via Layer-1 (type matching) + Layer-2 (LLM semantic judgment)
    \item Execute Z3 consistency pre-check; inject passing paradigms as hints
    \item Generate and verify LLM inference steps via Z3
    \item Upon UNSAT, extract UNSAT core and perform causal attribution
    \item Record typed repairs in Repair Trajectory Memory $\mathcal{M}$
    \item Detect stagnation (Jaccard similarity) and loops (cosine similarity); trigger strategy escalation
    \item Write successful repairs back to augment $\mathcal{L}$
\end{enumerate}

\subsection{Offline Stage: Paradigm Distillation}

\subsubsection{Trajectory Collection}

Given $N = 600$ training puzzles across difficulty levels, we execute LLM+Z3 solving $r = 3$ times per puzzle with $\text{temperature} = 0.7$ for inference diversity. We retain only successful trajectories where Z3 returns SAT with assignments matching ground truth. This yields approximately 1500 successful trajectories as the distillation source.

Each trajectory step records: Z3 verification labels (SAT/UNSAT/ERROR), domain change records $\Delta_t$ (variables with reduced domains), constraint type signatures, and preliminary step classification.

\subsubsection{Key Decision Point Identification}

A step $(S_{t-1}, a_t, S_t)$ is a Key Decision Point if:
\begin{itemize}
    \item \textbf{Condition A (Significant Domain Reduction):} $\exists x \in \Delta_t: |\delta_t(x)| \leq 1$ (a variable's domain becomes singleton)
    \item \textbf{Condition B (Non-Trivial Step Type):} Step classification is CHAIN\_PROPAGATION or CONTRADICTION\_ELIMINATION
\end{itemize}

Step types are assigned via lightweight rule-based classification on natural language action text. Across 1500 trajectories, we extract approximately 6000 KDPs.

\subsubsection{Cross-Trajectory Clustering and Paradigm Abstraction}

\textbf{Feature Representation:} Each KDP is represented as:
$$\mathbf{f}(\text{kdp}) = [\mathbf{h}_{\text{ct}}, \mathbf{b}_{\text{dom}}, \mathbf{e}_\tau, n_{\text{var}}, n_{\text{con}}]$$

where $\mathbf{h}_{\text{ct}}$ is a constraint-type histogram, $\mathbf{b}_{\text{dom}}$ is domain-size distribution, $\mathbf{e}_\tau$ is step-type one-hot encoding, and normalized problem scale metrics.

\textbf{Clustering:} We apply complete-linkage agglomerative clustering with cosine distance and threshold $\theta = 0.25$ (determined via grid search on development set). Clusters with fewer than 5 members are discarded.

\textbf{Paradigm Abstraction:} For each cluster, we uniformly sample up to 10 representative KDPs, format their solver states and reasoning text, and prompt an LLM to synthesize a reusable paradigm. LLM outputs are parsed into structured (trigger, operation, precondition, postcondition, scope) tuples. Up to 2 retries per cluster are permitted.

\subsubsection{Z3 Triple Verification}

For each candidate paradigm, we execute three independent verification checks:

\textbf{(1) Soundness:} Sample 50 random puzzle states satisfying trigger conditions. Test whether paradigm operations maintain Z3 satisfiability:
$$\text{soundness}(P) = \frac{1}{50} \sum_{i=1}^{50} \mathbf{1}[\text{Z3-SAT}(C_i \cup \{\text{op}(P)\})] \geq 0.90$$

\textbf{(2) Effect:} Verify that paradigm application actually reduces domain uncertainty (increases variable assignments):
$$\text{effect}(P) = \frac{1}{N_e} \sum_{i=1}^{N_e} \mathbf{1}[|\text{assignments}_{\text{after}}| > |\text{assignments}_{\text{before}}|] \geq 0.80$$

\textbf{(3) Trigger Precision:} Sample 20 states NOT satisfying trigger conditions; verify operations yield UNSAT (no false triggering):
$$\text{precision}(P) = \frac{1}{20} \sum_{i=1}^{20} \mathbf{1}[\text{Z3-UNSAT}(C_i \cup \{\text{op}(P)\})] \geq 0.80$$

Paradigms passing all three checks are persisted in Paradigm Library $\mathcal{L}$. Across 600 training puzzles, offline distillation yields approximately 35 verified paradigms (40:1 compression ratio).

\subsection{Online Stage: Paradigm-Guided Inference}

\subsubsection{Two-Layer Paradigm Retrieval}

For puzzle $\mathcal{P}^*$, at each step $t$, we extract constraint-type signature $\sigma(S_t)$ and apply:

\textbf{Layer-1 (Type Matching):} Return paradigms $P \in \mathcal{L}$ where required types $\subseteq \sigma(S_t)$ and confidence $\geq 0.50$. This typically reduces candidates from 35 to $< 5$.

\textbf{Layer-2 (Semantic Filtering):} When Layer-1 yields $\geq 2$ candidates, query LLM: "Given current solver state, which paradigm trigger conditions are genuinely satisfied?" Return top-3 candidates by confidence.

\subsubsection{Consistency Pre-Check and Hint Injection}

Before paradigm injection, execute Z3 consistency verification:
$$\text{consistent}(P, S_t) = \text{Z3-SAT}(C_t \cup \{\text{op}(P)\})$$

Only paradigms returning SAT are injected as "soft" (advisory) hints, not hard constraints. This prevents LLM misapplication when paradigms don't fully align with current state.

\subsubsection{Real-Time Z3 Verification and UNSAT Attribution}

After LLM generates inference action $a_t$, add to solver via constraint registry (enabling tracking) and call Z3:

\textbf{SAT Case:} Update domain assignments, save as potential revert checkpoint. Continue solving.

\textbf{UNSAT Case:} Extract UNSAT core $U_t$ and perform causal attribution:
$$\text{attr}(a_t, U_t) = \begin{cases} \text{NEW\_ASSERTION} & \text{if } a_t \in U_t \\ \text{LEGACY\_ERROR} & \text{if } a_t \notin U_t \end{cases}$$

If LEGACY\_ERROR, contradiction predates current action; revert $a_t$ and enter repair. If NEW\_ASSERTION, current action introduced contradiction; proceed to repair with UNSAT core as diagnostic.

\subsection{Repair Trajectory Memory}

\subsubsection{Memory Data Structure}

Repair Trajectory Memory $\mathcal{M}$ maintains ordered records $\{R_1, \ldots, R_L\}$, each a six-tuple:
$$R_t = (e_t, U_t, \hat{U}_t, \alpha_t, o_t, U'_t)$$

where:
\begin{itemize}
    \item $e_t$ is error type classification
    \item $U_t$ is UNSAT core (constraint IDs)
    \item $\hat{U}_t = \text{SHA256}(\text{sorted}(U_t))$ is core fingerprint for rapid comparison
    \item $\alpha_t$ is repair action (type, target constraint, summary text, embedding)
    \item $o_t$ is repair outcome (SAT or UNSAT)
    \item $U'_t$ is new UNSAT core post-repair (if UNSAT)
\end{itemize}

Repair action embeddings are computed via all-MiniLM-L6-v2 sentence transformer (384-dim), enabling semantic similarity queries.

\subsubsection{Stagnation Detection}

Stagnation occurs when recent repair records show high UNSAT core overlap (implying repairs are ineffective at addressing root causes):
$$\text{stagnate}(\mathcal{M}, k=3) = \mathbf{1}[\min_{i,j} J(U_i, U_j) \geq 0.75]$$

where Jaccard similarity is $J(A,B) = |A \cap B| / |A \cup B|$ computed across the last 3 records.

\subsubsection{Loop Detection}

Loop detection identifies when proposed repairs are semantically similar to historical attempts:
$$\text{loop}(\alpha_{\text{new}}) = \mathbf{1}[\max_t \cos(\mathbf{e}_{\text{new}}, \mathbf{e}_t) \geq 0.90]$$

\subsubsection{Four-Level Strategy Escalation}

Upon stagnation or loop detection:
\begin{enumerate}
    \item \textbf{L1 - Switch Target:} Preserve repair type but target different constraint from UNSAT core
    \item \textbf{L2 - Switch Type:} Change repair strategy class (e.g., relax-bound → split-constraint)
    \item \textbf{L3 - Revert Checkpoint:} Roll back to most recent SAT state; resume from alternative inference path
    \item \textbf{L4 - Full Retranslation:} Discard current formalization; re-translate all constraints from natural language with error history as guidance
\end{enumerate}

\section{Experiments}

\subsection{Experimental Setup}

\textbf{Benchmark:} ZebraLogic (Lin et al., 2025) contains logical grid puzzles ranging from 3×5 to 6×6 (30 to 180 variables, 15-80 constraints). We evaluate on three representative sizes: 4×4 (easy), 5×5 (medium), 6×6 (hard).

\textbf{Baselines:}
\begin{enumerate}
    \item Chain-of-Thought (CoT) prompting
    \item Logic-LM (Pan et al., 2023)
    \item ExpeL with CSP adaptation
    \item Neural-Symbolic baseline (LLM+Z3, no memory)
    \item Verified Few-Shot (direct paradigm injection without library)
    \item PRISM ablations: (a) memory only, (b) paradigm library only
\end{enumerate}

\textbf{Metrics:} Solving accuracy, average LLM calls per puzzle, average repair rounds, paradigm trigger rate, paradigm hit rate.

\subsection{Results}

[This section will contain: (1) Accuracy improvements, (2) Cost analysis (LLM calls), (3) Repair efficiency (stagnation reduction), (4) Paradigm transfer analysis, (5) Ablation studies]

\section{Conclusion}

PRISM introduces a novel paradigm for memory-augmented neural-symbolic reasoning specifically designed for CSP solving. By leveraging Z3 as a quality oracle for experience, encoding formally-verifiable paradigms, and implementing adaptive repair strategies via typed trajectory memory, PRISM achieves consistent improvements over neural-symbolic and memory-based baselines.

Key findings: (1) experience accumulation becomes exponentially valuable for harder problems, (2) formal verification of paradigms enables reliable knowledge transfer, (3) adaptive strategy escalation effectively breaks repair stagnation, (4) paradigms transfer across problem scales and types.

Future work includes: extending to other symbolic domains beyond CSPs, investigating curriculum learning for paradigm discovery, and exploring interaction with classical CSP solvers.

\section*{Acknowledgments}

We acknowledge support from [funding sources]. We thank [individuals] for valuable discussions.

\bibliographystyle{aaai2026}
\bibliography{aaai2026}

\end{document}
